% ==========================================
% ФАЙЛ: p4.tex
% БЛОК БИЛЕТОВ 16-20 (Определенный интеграл)
% ==========================================

% --- Вопрос 16 ---
\examquestion{16 (lucky)}{Свойства определенного интеграла (линейность и аддитивность)}
{Сформулируйте основные свойства определенного интеграла Римана. Докажите свойства линейности и аддитивности интеграла по промежутку интегрирования.}

\paragraph{Суть на пальцах (Краткая памятка)}
Интеграл — это, грубо говоря, площадь. Из этого логично вытекает:
\begin{itemize}
    \item \textbf{Линейность:} Если растянуть график в $k$ раз, площадь тоже вырастет в $k$ раз (константу выносим). Площадь суммы двух фигур равна сумме их площадей.
    \item \textbf{Аддитивность:} Если разрезать фигуру вертикальной линией на две части, то общая площадь равна сумме площадей левой и правой частей: $\int_a^b = \int_a^c + \int_c^b$.
\end{itemize}

\paragraph{Строгие формулировки}
Пусть функции $f(x)$ и $g(x)$ интегрируемы на отрезке $[a, b]$ (обозначается $f, g \in \mathcal{R}[a, b]$).
\begin{enumerate}
    \item \textbf{Линейность:} Для любых констант $\alpha, \beta \in \mathbb{R}$ функция $(\alpha f(x) + \beta g(x))$ также интегрируема на $[a, b]$, причем:
    $$ \int_a^b (\alpha f(x) + \beta g(x)) dx = \alpha \int_a^b f(x) dx + \beta \int_a^b g(x) dx $$
    \item \textbf{Аддитивность:} Для любой точки $c \in (a, b)$ функция интегрируема на $[a, c]$ и $[c, b]$, причем:
    $$ \int_a^b f(x) dx = \int_a^c f(x) dx + \int_c^b f(x) dx $$
\end{enumerate}

\paragraph{Доказательство линейности}
Основано на свойствах конечных сумм. Запишем интегральную сумму для линейной комбинации:
$$ \sigma(\alpha f + \beta g, T, \xi) = \sum_{k=1}^n \big(\alpha f(\xi_k) + \beta g(\xi_k)\big) \Delta x_k = \alpha \sum_{k=1}^n f(\xi_k)\Delta x_k + \beta \sum_{k=1}^n g(\xi_k)\Delta x_k $$
Переходя к пределу при мелкости разбиения $\delta(T) \to 0$, пределы правых сумм существуют (т.к. $f, g \in \mathcal{R}$), следовательно, предел левой части тоже существует и равен сумме интегралов.

\paragraph{Доказательство аддитивности}
\begin{itemize}
    \item Пусть предел $\int_a^b$ существует. Так как предел интегральных сумм не зависит от выбора разбиения $T$, мы можем рассматривать \textbf{только такие разбиения}, которые обязательно включают точку $c$ в качестве одной из точек деления ($x_m = c$).
    \item Тогда любая интегральная сумма на $[a, b]$ жестко распадается на две независимые суммы:
    $$ \sum_{[a, b]} f(\xi_k)\Delta x_k = \sum_{[a, c]} f(\xi_k)\Delta x_k + \sum_{[c, b]} f(\xi_k)\Delta x_k $$
    \item Переходя к пределу при $\delta(T) \to 0$, предел левой части существует по условию, следовательно, существуют и пределы слагаемых (интегралы по частям отрезка), и их сумма равна исходному интегралу.
\end{itemize}


% --- Вопрос 17 ---
\examquestion{17}{Первая теорема о среднем}
{Сформулируйте и докажите первую теорему о среднем значении для определенного интеграла. Дайте геометрическую интерпретацию.}

\paragraph{(Суть на пальцах (Краткая памятка))}
Представь кривую гору песка (наш график). Если этот песок разровнять бульдозером, получится плоский прямоугольник той же ширины и такой же площади. 
Высота этого ровного прямоугольника — это и есть \textbf{среднее значение функции}. Теорема гласит: если функция непрерывна (без разрывов), то на графике обязательно найдется хотя бы одна точка $c$, в которой реальная высота горы совпадает с этой «усредненной» высотой.

\paragraph{Формулировка теоремы}
Если функция $f(x)$ \textbf{непрерывна} на отрезке $[a, b]$, то на этом отрезке существует хотя бы одна точка $c \in [a, b]$, такая что:
$$ \int_a^b f(x) dx = f(c) \cdot (b - a) $$
Число $f(c) = \frac{1}{b-a} \int_a^b f(x) dx$ называется средним значением функции на отрезке.

\paragraph{Доказательство (Строго по шагам)}
\begin{enumerate}
    \item Так как $f(x)$ непрерывна на замкнутом отрезке $[a, b]$, по теореме Вейерштрасса она достигает на нем своего минимума $m$ и максимума $M$. То есть: $\forall x \in [a, b] \implies m \le f(x) \le M$.
    \item Проинтегрируем это двойное неравенство на отрезке $[a, b]$ (используем свойство сохранения знака и линейность интеграла):
    $$ \int_a^b m \, dx \le \int_a^b f(x) dx \le \int_a^b M \, dx $$
    \item Вычислим крайние интегралы, Эм большое и маленькое тут константы, вынесем. (константы выносятся, остается интеграл от $dx$, равный длине отрезка):
    $$ m(b - a) \le \int_a^b f(x) dx \le M(b - a) $$
    \item Разделим всё на положительную длину отрезка $(b - a)$:
    $$ m \le \underbrace{\frac{1}{b - a} \int_a^b f(x) dx}_{\text{Обозначим это число за } \mu} \le M $$
    \item Мы получили число $\mu$, которое лежит между минимумом $m$ и максимумом $M$ функции. По теореме Коши (о промежуточном значении непрерывной функции), функция обязана принимать любое значение между своим минимумом и максимумом. Значит, найдется точка $c \in [a, b]$, где $f(c) = \mu$.
    \item Подставляя $\mu = f(c)$ обратно в дробь из пункта 4, получаем: $\int_a^b f(x) dx = f(c)(b - a)$. Теорема доказана.
\end{enumerate}


% --- Вопрос 18 ---
\examquestion{18}{Интеграл с переменным верхним пределом}
{Введите понятие интеграла с переменным верхним пределом. Сформулируйте и докажите теорему Барроу (о производной интеграла по верхнему пределу). Какова роль этой теоремы в матанализе?}

\paragraph{Суть на пальцах (Краткая памятка)}
Мы привыкли, что $\int_a^b$ — это конкретное число (площадь). А что если мы зафиксируем только старт $a$, а правую границу $x$ будем двигать вправо-влево? Площадь будет меняться в зависимости от $x$. Получается \textbf{функция площади} $\Phi(x) = \int_a^x f(t)dt$. 
Букву под интегралом поменяли на $t$, чтобы не путать её с движущейся границей $x$.
\textbf{Теорема Барроу:} Скорость, с которой прирастает площадь ($\Phi'(x)$), в точности равна высоте графика в этой точке ($f(x)$). 

\paragraph{Формулировка (Теорема Барроу)}
Пусть $f(t)$ интегрируема на $[a, b]$. Рассмотрим функцию верхнего предела $\Phi(x) = \int_a^x f(t) dt$.
Если подынтегральная функция $f(t)$ \textbf{непрерывна} в точке $x$, то функция $\Phi(x)$ дифференцируема в этой точке, причем её производная равна:
$$ \Phi'(x) = f(x) $$
\textit{Роль теоремы: Она доказывает, что у любой непрерывной функции гарантированно существует первообразная (и это функция $\Phi(x)$).}

\paragraph{Доказательство}
Докажем по определению производной: $\Phi'(x) = \lim_{\Delta x \to 0} \frac{\Phi(x+\Delta x) - \Phi(x)}{\Delta x}$.
\begin{enumerate}
    \item Запишем приращение функции площади $\Delta \Phi$:
    $$ \Delta \Phi = \Phi(x+\Delta x) - \Phi(x) = \int_a^{x+\Delta x} f(t)dt - \int_a^x f(t)dt $$
    \item По свойству аддитивности, разность этих площадей — это просто площадь маленького довеска от $x$ до $x+\Delta x$:
    $$ \Delta \Phi = \int_x^{x+\Delta x} f(t)dt $$
    \item Применим к этому маленькому интегралу Первую теорему о среднем (Вопрос 17). Найдется точка $c$, лежащая между $x$ и $x+\Delta x$, такая что:
    $$ \Delta \Phi = f(c) \cdot \big((x+\Delta x) - x\big) = f(c) \cdot \Delta x $$
    \item Составим отношение приращения функции к приращению аргумента:
    $$ \frac{\Delta \Phi}{\Delta x} = f(c) $$
    \item Перейдем к пределу при $\Delta x \to 0$. Так как точка $c$ зажата между $x$ и $x+\Delta x$, по теореме о двух милиционерах точка $c \to x$.
    \item А так как $f$ непрерывна в точке $x$, то $\lim_{c \to x} f(c) = f(x)$. Следовательно:
    $$ \Phi'(x) = \lim_{\Delta x \to 0} \frac{\Delta \Phi}{\Delta x} = \lim_{c \to x} f(c) = f(x) $$
\end{enumerate}


% --- Вопрос 19 ---
\examquestion{19}{Формула Ньютона-Лейбница}
{Сформулируйте и докажите формулу Ньютона-Лейбница, связывающую определенный интеграл с первообразной. Приведите пример вычисления.}

\paragraph{Суть на пальцах (Краткая памятка)}
Это \textbf{главная теорема в матанализе. Это БАЗА.}. Она связывает две абсолютно разные вещи: поиск первообразной (действие, обратное производной) и поиск площади под кривой (предел дробных сумм прямоугольников).
Формула говорит: не нужно считать бесконечные пределы сумм! Найди первообразную $F(x)$, подставь верхнюю границу, вычти нижнюю границу — и площадь готова.

\paragraph{Формулировка}
Если функция $f(x)$ непрерывна на отрезке $[a, b]$, а $F(x)$ — любая её первообразная на этом отрезке (т.е. $F'(x) = f(x)$), то справедлива формула:
$$ \int_a^b f(x) dx = F(b) - F(a) \equiv F(x)\Big|_a^b $$

\paragraph{Доказательство}
\begin{enumerate}
    \item Из теоремы Барроу (Вопрос 18) мы знаем, что интеграл с переменным верхним пределом $\Phi(x) = \int_a^x f(t) dt$ является первообразной для $f(x)$, так как $\Phi'(x) = f(x)$.
    \item По теореме о первообразных, любые две первообразные одной и той же функции отличаются только на константу $C$. Следовательно, наша произвольная первообразная $F(x)$ связана с $\Phi(x)$ уравнением:
    $$ F(x) = \Phi(x) + C = \int_a^x f(t)dt + C $$
    \item Чтобы найти константу $C$, подставим $x = a$. Интеграл от $a$ до $a$ равен нулю (площадь отрезка нулевой ширины):
    $$ F(a) = \int_a^a f(t)dt + C = 0 + C \implies C = F(a) $$
    \item Теперь перепишем уравнение из пункта 2, зная $C$, и подставим верхнюю границу $x = b$:
    $$ F(b) = \int_a^b f(t)dt + F(a) $$
    \item Переносим $F(a)$ в левую часть и получаем искомую формулу:
    $$ \int_a^b f(x) dx = F(b) - F(a) $$
\end{enumerate}





% --- Вопрос 20 ---
\examquestion{20}{Замена переменной и интегрирование по частям}
{Сформулируйте и докажите теоремы о замене переменной и формуле интегрирования по частям для определенного интеграла Римана. Укажите главное отличие от неопределенного интеграла.}

\paragraph{Главное отличие от неопределенного интеграла}
\begin{itemize}
    \item \textbf{Замена переменной ($x \to t$):} Пределы интегрирования $a, b$ \textbf{обязательно пересчитываются} в новые числа $\alpha, \beta$. Возвращаться к старой переменной $x$ в конце не нужно.
    \item \textbf{Интегрирование по частям:} Переменная не меняется, пределы те же, но внеинтегральный член сразу вычисляется на границах: $\Big[ U(x)V(x) \Big]_a^b$.
\end{itemize}

\paragraph{Теорема 1 (Замена переменной)}
Пусть $f(x) \in \mathcal{C}[a, b]$, а $x = \varphi(t)$ непрерывно дифференцируема на $[\alpha, \beta]$, причем $\varphi(\alpha) = a$, $\varphi(\beta) = b$ и $\varphi([\alpha, \beta]) \subseteq [a, b]$. Тогда:
$$ \int_a^b f(x) dx = \int_\alpha^\beta f(\varphi(t)) \cdot \varphi'(t) dt $$

\textit{Идея доказательства:} Дифференцирование сложной функции + формула Ньютона-Лейбница (ФНЛ).
\begin{proof}
Пусть $F(x)$ — первообразная для $f(x)$, то есть $F'(x) = f(x)$. По правилу сложной функции:
$$ (F(\varphi(t)))'_t = F'(\varphi(t)) \cdot \varphi'(t) = f(\varphi(t)) \cdot \varphi'(t) $$
Значит, $F(\varphi(t))$ — первообразная для правой части. Применим ФНЛ к обоим интегралам:
\begin{itemize}
    \item Левый: $\int_a^b f(x) dx = F(b) - F(a)$
    \item Правый: $\int_\alpha^\beta f(\varphi(t))\varphi'(t) dt = F(\varphi(t))\Big|_\alpha^\beta = F(\varphi(\beta)) - F(\varphi(\alpha)) = F(b) - F(a)$
\end{itemize}
Левые части равны, так как равны правые.
\end{proof}

\paragraph{Теорема 2 (Интегрирование по частям)}
Пусть $U(x), V(x)$ непрерывно дифференцируемы на $[a, b]$. Тогда:
$$ \int_a^b U(x) V'(x) dx = \Big[ U(x)V(x) \Big]_a^b - \int_a^b V(x) U'(x) dx $$

\textit{Идея доказательства:} Интегрирование производной произведения функций + ФНЛ.
\begin{proof}
По правилу Лейбница: $(UV)' = U'V + UV'$. Проинтегрируем равенство на $[a, b]$:
$$ \int_a^b (UV)' dx = \int_a^b U'V dx + \int_a^b UV' dx $$
По ФНЛ левый интеграл равен $\int_a^b (UV)' dx = [UV]_a^b$. Подставим и выразим нужный интеграл:
$$ \Big[ U(x)V(x) \Big]_a^b = \int_a^b V dU + \int_a^b U dV \implies \int_a^b U dV = \Big[ U(x)V(x) \Big]_a^b - \int_a^b V dU $$
\end{proof}

\paragraph{Бонус-пример (Комбо: Замена с пересчетом пределов $\to$ По частям)}
Вычислим $I = \int_0^{\pi^2/4} \cos(\sqrt{x}) dx$.

\textbf{1. Замена переменной:} 
Пусть $t = \sqrt{x} \implies x = t^2 \implies dx = 2t \, dt$.
\textit{Пересчет пределов:} 
\begin{itemize}
    \item Нижний: $x = 0 \implies \mathbf{t = 0} \equiv \alpha$
    \item Верхний: $x = \pi^2/4 \implies \mathbf{t = \pi/2} \equiv \beta$
\end{itemize}
Получаем новый интеграл (работаем только в пределах от $0$ до $\pi/2$ по переменной $t$):
$$ I = \int_0^{\pi/2} \cos(t) \cdot 2t \, dt = 2\int_0^{\pi/2} t \cos t \, dt $$

\textbf{2. Интегрирование по частям (в новых пределах $\alpha, \beta$):}
Пусть $U = t \implies dU = dt$ \\
Пусть $dV = \cos t \, dt \implies V = \sin t$.
$$ \int_0^{\pi/2} t \cos t \, dt = \Big[ t \sin t \Big]_0^{\pi/2} - \int_0^{\pi/2} \sin t \, dt $$
\begin{itemize}
    \item Вычисляем внеинтегральный член на новых пределах: $\Big[ t \sin t \Big]_0^{\pi/2} = \left(\frac{\pi}{2} \cdot \sin\frac{\pi}{2}\right) - (0 \cdot \sin 0) = \frac{\pi}{2} - 0 = \frac{\pi}{2}$
    \item Добиваем оставшийся интеграл: $\int_0^{\pi/2} \sin t \, dt = \Big[ -\cos t \Big]_0^{\pi/2} = \left(-\cos\frac{\pi}{2}\right) - (-\cos 0) = 0 + 1 = 1$
\end{itemize}
Собираем всё вместе, не забывая внешнюю двойку:
$$ I = 2 \cdot \left( \frac{\pi}{2} - 1 \right) = \pi - 2 $$